% =============================================================================
% Project : THE_PHILOSOPHY_OF_VIABILITY
% Book    : Book I — Ontology of Viability
% Book ID : VS.B01
%
% File    : VS.B01_ONTOLOGY_OF_VIABILITY/2_MAINMATTER/PART_I_IDENTITY_AND_MANIFESTATION/CHAPTER_07_FORMAL_AND_LOGICAL_RECONSTRUCTION/sections/section_16_claim_status_and_justificatory_classification.tex
%
% Section / Chapter:
%   Chapter 07, Section 16
%   Claim Status and Justificatory Classification
%
% Purpose :
%   To establish a two-axis classification system for claims,
%   distinguishing their functional or logical form from their
%   epistemic and justificatory status; to clarify the conditions
%   under which claims may be adopted, stipulated, argued,
%   conditionally supported, treated as provisional, retained
%   as hypotheses, classified as open problems, or rejected.
%
% =============================================================================

\section{Claim Status and Justificatory Classification}
\label{sec:claim-status-and-justificatory-classification}

\subsection{The Need for a Two-Axis Classification}

A philosophical system contains claims that perform
different functions and possess different degrees
or kinds of justification.

Some claims establish the commitments from which
a particular inquiry proceeds. Others stipulate
the meanings of technical terms, formulate
propositions, present arguments, introduce
hypotheses, or identify questions that remain
unresolved.

These differences cannot be represented adequately
by a single classification.

A statement may function as a definition while
its underlying conceptual choice remains provisional.
A proposition may be formally derived within a
specified system while the premises from which
it is derived remain conditionally supported.
A foundational commitment may govern the construction
of an ontology without thereby being an empirically
verified or deductively demonstrated conclusion.

The classification of claims must therefore
distinguish two questions:

\begin{enumerate}
    \item What is the functional or logical form
    of the claim?
    \item What is the epistemic or justificatory
    status of the claim?
\end{enumerate}

The first question concerns the role a claim
performs within the philosophical system.

The second concerns the grounds on which the
claim is adopted, maintained, supported,
questioned, or rejected.

The central thesis of this section is that
these two dimensions are independent axes
of classification. Neither can replace
the other, and neither alone provides
a complete assessment of a claim.

This distinction is necessary for the
philosophical reconstruction of Book I.
It allows foundational commitments to be
identified without treating them as
self-proving, definitions to be used
without confusing stipulation with truth,
and derived propositions to be assessed
without transferring the certainty of
a valid inference to premises that have
not been independently established.

\subsection{The First Axis: Functional and Logical Form}

The first axis classifies claims according
to their role within the structure of
philosophical inquiry.

The principal forms used in this book are:

\begin{enumerate}
    \item Foundational Commitment;
    \item Axiom;
    \item Definition;
    \item Proposition;
    \item Thesis;
    \item Hypothesis;
    \item Argument;
    \item Open Problem.
\end{enumerate}

These categories identify different
functions. They do not constitute
a ranking from weakest to strongest.

\subsubsection{Foundational Commitment}

A Foundational Commitment is a principle
or constraint adopted as a governing
orientation of the philosophical
system.

It determines what the system is
committed to preserving, examining,
or refusing to assume in the course
of its construction.

A foundational commitment may
establish the conceptual and
ontological conditions under
which subsequent arguments
are developed.

For example, the three principles
developed in the examination of
Being and Identity are:

\begin{itemize}
    \item Analytical Distinction;
    \item Ontological Inseparability;
    \item Absence of Ontological Priority.
\end{itemize}

Their classification as foundational
commitments identifies their role
within the ontology.

It does not mean that one principle
automatically proves the others,
nor that their designation
eliminates the need for philosophical
argument.

A foundational commitment may
receive justification through
conceptual analysis, argument,
critical examination, or
its demonstrated necessity
for the coherent development
of the system.

Its foundational role must
not be confused with an
unquestionable epistemic status.

\subsubsection{Axiom}

An Axiom is a statement adopted
as a starting assumption within
a specified formal system.

Its function is formal.
It establishes a premise from
which further propositions
may be derived according
to the rules of the relevant
calculus.

An axiom may formalize
a foundational commitment,
but the two categories
are not identical.

A foundational commitment
belongs to the philosophical
architecture of the system.
An axiom belongs to a
specified formal formulation.

The transition from one
to the other requires
an explicit formalization
and an account of the
interpretation under which
the axiom is adopted.

The designation of a statement
as an axiom does not establish
its ontological truth.

It establishes its role
within a formal system.

\subsubsection{Definition}

A Definition specifies
the meaning or formal
conditions of application
of a term, symbol, or
expression within a
specified context.

Definitions may be
stipulative, formal,
operational, or otherwise
appropriate to the
relevant inquiry.

A definition establishes
how an expression is
to be used within the
specified framework.

It does not, by itself,
prove that the phenomenon
being described exists
or that the defined
concept accurately
captures its ontological
subject.

For example, defining
a formal continuity
relation by reference
to stages, respects,
and criteria establishes
a formal convention
for examining continuity.
It does not prove
that every actual
continuity relation
satisfies that convention.

A definition must therefore
be assessed according
to clarity, consistency,
scope, conceptual adequacy,
and compatibility with
the ontological commitments
of the system.

\subsubsection{Proposition}

A Proposition is a claim
presented for examination
within a specified
conceptual or formal
framework.

It may be established
through argument, derived
from adopted premises,
or proposed for critical
assessment.

A proposition is not
necessarily a theorem.

Its status as a proposition
identifies its role as
a claim whose content
and support can be
examined.

Where a proposition
is formally derived,
the relevant premises
and inference rules
must be stated.

Where it is supported
through philosophical
argument, the reasoning
must be made explicit.

The word proposition
therefore identifies
a logical form and
not a degree of certainty.

\subsubsection{Thesis}

A Thesis is a substantive
philosophical position
advanced for defense
through a structured
body of reasons.

A thesis generally
organizes several claims,
arguments, distinctions,
or interpretations
around a central
philosophical assertion.

It may be supported
to different degrees
and may contain
elements whose
justificatory statuses
differ.

A thesis should not
be treated as a single
indivisible assertion
when its components
have distinct logical
or justificatory roles.

Its evaluation requires
the examination of
its central claim,
supporting arguments,
presuppositions,
and consequences.

\subsubsection{Hypothesis}

A Hypothesis is a proposed
claim or explanatory
construction introduced
for examination when
its status has not
been sufficiently
established.

A hypothesis may
organize inquiry,
generate testable
or analytically
examinable consequences,
or provide a possible
interpretation of
a phenomenon.

Its role is exploratory
or provisional in
relation to the
question under
investigation.

A hypothesis is
not equivalent to
an arbitrary speculation.
It must be formulated
clearly enough to
permit examination,
comparison, or
critical assessment.

However, the usefulness
of a hypothesis for
organizing inquiry
does not establish
its truth.

\subsubsection{Argument}

An Argument is a structured
set of premises and
inferential relations
presented in support
of a conclusion.

Its formal role
is to show how
a conclusion is
supported by
specified premises.

The classification
of a passage as
an argument does
not determine whether
the argument is
sound, valid,
cogent, or adequate.

Those questions
require independent
evaluation of its
premises, inference
structure, and
interpretation.

An argument may
be formally valid
while containing
premises that are
not adequately
justified.

It may also be
philosophically
persuasive without
being represented
as a formal proof.

The distinction
between argumentative
form and justificatory
success must therefore
be preserved.

\subsubsection{Open Problem}

An Open Problem
is a question,
difficulty, or
unresolved relation
that has not been
satisfactorily settled
within the current
development of the
system.

It may arise from
incomplete definitions,
insufficient premises,
competing interpretations,
unresolved consequences,
or limitations of
the available methods.

An open problem
is not necessarily
a defect in the
system.

Its explicit recognition
may be a condition
of philosophical
precision, particularly
where a premature
answer would impose
an unjustified
ontological commitment.

Open problems
must nevertheless
be formulated
with sufficient
specificity to
distinguish them
from general
statements of
uncertainty.

\subsection{The Second Axis: Epistemic and Justificatory Status}

The second axis concerns
the grounds on which
a claim is maintained
within the system.

The principal statuses
are:

\begin{enumerate}
    \item Foundationally Adopted Commitments;
    \item Stipulated or Defined Claims;
    \item Argued or Derived Claims;
    \item Conditionally Supported Claims;
    \item Provisional Constructions;
    \item Hypotheses;
    \item Unresolved or Open Problems;
    \item Rejected or Superseded Claims.
\end{enumerate}

These categories identify
the justificatory standing
of a claim within
the present development
of the book.

They do not necessarily
describe permanent
properties of claims.

A claim may acquire
stronger support,
be reformulated,
be restricted to
a narrower scope,
or be rejected
as the inquiry
develops.

The status assigned
to a claim must
therefore be traceable
to its actual grounds
and current role
within the system.

\subsubsection{Foundationally Adopted Commitments}

This status applies
to commitments adopted
as governing principles
of the present
philosophical construction.

They provide the
starting orientation
for the development
of the ontology.

Their adoption
must be explicit.

A foundationally
adopted commitment
is not thereby
a proven theorem
or an empirically
verified conclusion.

Its justification
may be developed
through the philosophical
arguments of the
book, but its
current function
must remain distinguishable
from the arguments
offered in its support.

The three principles
concerning the
analytical distinction,
ontological inseparability,
and absence of
ontological priority
between Being and
Identity provide
examples of commitments
whose roles must
be identified
without conflating
their individual
justifications.

\subsubsection{Stipulated or Defined Claims}

This status applies
to claims whose
content is established
through an explicit
definition or
stipulation within
a specified framework.

Their immediate
justification is
the declared
convention governing
the use of the
relevant expression.

A stipulated
definition may
be evaluated
for coherence,
precision, adequacy,
and compatibility
with the wider
ontology.

Its status does
not establish
that the defined
term corresponds
to an independently
existing ontological
entity or relation.

Where a definition
contains substantive
ontological assumptions,
those assumptions
must be identified
and evaluated
separately.

\subsubsection{Argued or Derived Claims}

This status applies
to claims supported
through explicit
reasoning or
derived from
specified premises.

A claim may
be argued through
conceptual analysis,
philosophical reasoning,
comparison, or
other relevant
methods.

A claim may
be derived formally
through a sequence
of valid inference
steps.

These forms
of support must
be distinguished.

For a formally
derived claim,
the premises,
formal language,
and inference
rules must be
identifiable.

For an argued
claim, the
supporting reasons
and the relation
between those
reasons and
the conclusion
must be examinable.

The designation
does not guarantee
that every argument
is sound or
that every derivation
begins from
ontologically
established premises.

It records the
presence of
an articulated
justificatory basis.

\subsubsection{Conditionally Supported Claims}

This status applies
to claims supported
under specified
assumptions, conditions,
interpretations, or
restricted domains.

A conditionally
supported claim
may be well
justified within
its stated scope
without being
established as
universally applicable.

Its conditions
must be made
explicit.

For example,
a formal implication
may hold within
a particular
calculus or
under a specified
interpretation.

Its validity
within that
framework does
not automatically
establish its
unrestricted
ontological validity.

Conditional support
is therefore
not necessarily
weak support.

It may be
strong within
a precisely
defined domain
while remaining
limited in
scope.

\subsubsection{Provisional Constructions}

This status applies
to concepts, distinctions,
models, classifications,
or formal arrangements
introduced as
working constructions
whose adequacy
has not yet
been fully established.

A provisional
construction may
be useful for
organizing inquiry
and making
relationships
explicit.

Its provisional
status must
remain visible
so that the
construction is
not silently
converted into
a foundational
commitment.

A construction
may later be
refined, replaced,
restricted, or
rejected.

Its usefulness
as a working
instrument does
not establish
its ontological
necessity.

\subsubsection{Hypotheses}

This status applies
to claims proposed
for examination
whose support
is not yet
sufficient to
justify their
adoption as
established positions
within the system.

A hypothesis
must be sufficiently
clear to permit
critical assessment.

It may be
supported by
preliminary arguments
or explanatory
considerations,
but those grounds
must not be
represented as
conclusive when
they are not.

A hypothesis
may remain
open, acquire
conditional support,
become an argued
claim, or be
rejected as
the investigation
progresses.

\subsubsection{Unresolved or Open Problems}

This status applies
to questions or
claims for which
the available
analysis has
not established
a satisfactory
answer.

An unresolved
problem may
concern the
meaning of
a concept,
the relation
between two
concepts,
the validity
of a proposed
inference,
or the ontological
implications of
a formal construction.

Its classification
must identify
what remains
unresolved.

A problem
should not
be marked
as settled
merely because
a possible
answer has
been proposed.

The distinction
between a proposed
solution and
an established
resolution is
essential.

\subsubsection{Rejected or Superseded Claims}

This status applies
to claims that
are no longer
maintained within
the current
development of
the system.

A claim may
be rejected
because it
is contradictory,
circular, unsupported,
conceptually inadequate,
or incompatible
with commitments
that have been
retained.

A claim may
also be superseded
when a more
adequate formulation
replaces it.

Rejection and
supersession should
not be treated
as identical.

Rejection records
a negative
assessment of
the claim or
its justification.

Supersession records
that the claim
has been replaced
or reformulated,
which may occur
without every
part of its
earlier content
being rejected.

The reason
for rejection
or supersession
should be recorded
whenever the
claim has played
a significant
role in the
development
of the book.

\subsection{Independence of the Two Axes}

The two classification
axes answer different
questions.

Functional or
logical form
identifies what
a claim does
within the
system.

Epistemic or
justificatory status
identifies the
grounds on which
the claim is
maintained.

Consequently,
the same functional
form may occur
under different
justificatory statuses.

A proposition
may be argued
or derived,
conditionally supported,
provisional, or
unresolved.

A definition
may be stipulated
for formal use
while its broader
conceptual adequacy
remains provisional.

A thesis may
be argued
with substantial
support while
some of its
component claims
remain hypothetical.

An axiom may
be adopted
within a formal
calculus while
its ontological
interpretation
remains unresolved.

The relationship
between the
axes is therefore
many-to-many
rather than
a fixed correspondence.

Formally, let:

\[
    F
\]

denote the
set of functional
and logical
forms, and let:

\[
    J
\]

denote the
set of epistemic
and justificatory
statuses.

For a claim
\(c\), its
classification
may be represented
as:

\[
    \operatorname{Class}(c)
    =
    \langle
        f(c),j(c)
    \rangle,
\]

where:

\[
    f(c)\in F
\]

and:

\[
    j(c)\in J.
\]

This ordered
pair represents
two independent
classifications
of the same
claim.

It does not
imply that
every theoretically
possible combination
must occur.

Some combinations
may be inappropriate
or require
special explanation.

The important
point is that
the second
classification
must not be
inferred solely
from the first.

\subsection{Examples of Cross-Classification}

The following
examples illustrate
how the two
axes operate
together.

\paragraph{Example 1: A foundational commitment.}

A principle
may be classified
as:

\[
    \langle
        \text{Foundational Commitment},
        \text{Foundationally Adopted}
    \rangle.
\]

Its first
classification
identifies its
governing role.

Its second
records that
it has been
adopted as
a starting
commitment.

The pair
does not
mean that
the principle
has been
deductively
proved.

\paragraph{Example 2: A formal axiom.}

A statement
may be classified
as:

\[
    \langle
        \text{Axiom},
        \text{Conditionally Supported}
    \rangle.
\]

This indicates
that the statement
is adopted
as a formal
starting point
within a specified
system, while
its support
or interpretation
is limited
to stated
conditions.

Its axiomatic
role does
not remove
the conditions
governing its
justification.

\paragraph{Example 3: A definition.}

A formal
definition
may be classified
as:

\[
    \langle
        \text{Definition},
        \text{Stipulated or Defined}
    \rangle.
\]

The pair
records that
the expression
is being
used according
to an explicit
formal convention.

It does
not establish
that the
definition
is the
only possible
or ontologically
necessary
account.

\paragraph{Example 4: A proposition.}

A proposition
may be classified
as:

\[
    \langle
        \text{Proposition},
        \text{Argued or Derived}
    \rangle.
\]

Its status
records that
reasons or
inference steps
have been
provided in
its support.

The adequacy
of those
reasons must
still be
examined.

\paragraph{Example 5: A hypothesis.}

A proposed
account may
be classified
as:

\[
    \langle
        \text{Hypothesis},
        \text{Hypothesis}
    \rangle.
\]

The first
classification
identifies
its role
as a proposed
claim for
examination.

The second
records that
its justificatory
standing remains
hypothetical.

Although the
same term
appears on
both axes
in this example,
the two
classifications
perform different
functions.

\paragraph{Example 6: An unresolved problem.}

A question
may be classified
as:

\[
    \langle
        \text{Open Problem},
        \text{Unresolved or Open Problem}
    \rangle.
\]

The first
classification
identifies
its role
within the
structure of
inquiry.

The second
records that
no adequate
resolution
has yet
been established.

\subsection{The Difference Between Formal Derivation and Ontological Justification}

A particularly
important application
of the two-axis
system concerns
formal derivation.

Suppose a
conclusion \(q\)
is derived
from premises
\(p_1,\ldots,p_n\)
under a
formal calculus
\(\mathcal{L}\).

The derivation
may be represented
as:

\[
    p_1,\ldots,p_n
    \vdash_{\mathcal{L}} q.
\]

This notation
expresses formal
derivability
within the
specified calculus.

It does
not independently
establish the
ontological truth
of the premises.

The derivation
supports the
conclusion
relative to
the premises
and inference
rules.

Consequently,
the justificatory
status of
\(q\) must
be assessed
with reference
to the status
of the premises,
the adequacy
of the calculus,
and the
interpretation
connecting the
formal language
to the ontological
subject.

A formally
valid derivation
from conditionally
supported premises
does not
automatically
produce an
unconditionally
established
ontological claim.

Similarly,
a formal
theorem may
be established
within a
mathematical
structure without
thereby proving
that the
structure describes
the actual
constitution
of Being.

The distinction
between derivability
and ontological
justification
is therefore
essential.

\subsection{Claim Status and Conceptual Development}

Claims may
change their
justificatory
status during
the development
of a philosophical
system.

Such change
must be
recorded rather
than concealed.

A hypothesis
may receive
additional
support and
become a
conditionally
supported claim.

A provisional
construction
may be
refined into
a definition
or rejected
as conceptually
inadequate.

An argument
may reveal
that a proposition
depends upon
an unstated
premise.

A foundational
commitment
may receive
a more
developed
justification
without losing
its original
role.

These changes
do not imply
that the
functional form
and justificatory
status must
always change
together.

For example,
a proposition
may remain
a proposition
while its
justificatory
status changes
from provisional
to argued
or derived.

The classification
system must
therefore permit
the history
of a claim
to be recorded
without confusing
its present
status with
its earlier
formulations.

Where a
claim is
revised, the
relation between
the earlier
and later
versions should
be made
clear.

\subsection{Scope, Context, and Granularity}

The classification
of a claim
must be
sensitive to
its scope
and level
of articulation.

A complex
paragraph may
contain several
claims with
different functional
forms and
justificatory
statuses.

A thesis
may include
definitions,
arguments,
propositions,
and hypotheses.

A single
sentence may
contain both
a stipulated
condition and
a substantive
ontological assertion.

The classification
must therefore
be applied
at an appropriate
level of
granularity.

Where a
claim contains
several logically
distinct components,
those components
should be
separated for
classification
when their
statuses differ
materially.

The scope
of a claim
must also
be recorded.

A claim
supported for
one domain,
respect, or
interpretation
must not
be classified
as universally
supported unless
its wider
applicability
has been
established.

The same
requirement
applies to
claims concerning
Identity, Invariance,
Unity, Expression,
Manifestation,
Persistence, and
Continuity.

The classification
system must
preserve the
conditions under
which each
claim is
maintained.

\subsection{Criteria for Assigning Status}

The assignment
of a functional
form should
be based
on the
role of
the claim
within the
argumentative
and formal
structure.

The assignment
of justificatory
status should
be based
on the
actual grounds
available for
the claim.

The following
criteria guide
the classification.

\paragraph{Functional-form criteria.}

\begin{enumerate}
    \item Does the claim
    establish a governing
    commitment?
    \item Is it adopted
    as a starting
    statement in
    a specified
    formal calculus?
    \item Does it
    stipulate the
    meaning or
    use of
    an expression?
    \item Is it
    a claim
    presented for
    examination?
    \item Does it
    advance a
    substantive
    philosophical
    position?
    \item Is it
    proposed for
    exploratory
    examination?
    \item Does it
    contain premises
    and inferential
    support for
    a conclusion?
    \item Does it
    identify an
    unresolved
    question?
\end{enumerate}

\paragraph{Justificatory-status criteria.}

\begin{enumerate}
    \item Has the
    claim been
    explicitly adopted
    as a foundational
    commitment?
    \item Is its
    content established
    through stipulation
    or definition?
    \item Are explicit
    arguments or
    derivations available?
    \item Does its
    support depend
    upon stated
    conditions or
    restricted interpretations?
    \item Is it
    a working
    construction
    whose adequacy
    remains under
    examination?
    \item Is it
    a hypothesis
    awaiting sufficient
    justification?
    \item Does it
    remain unresolved?
    \item Has it
    been rejected
    or replaced
    by a
    later formulation?
\end{enumerate}

These criteria
do not replace
philosophical judgment.

They make
the grounds
of classification
explicit and
provide a
consistent basis
for reviewing
the claims
of the book.

\subsection{The Claim Register}

The two-axis
classification
may be implemented
through a claim
register.

A claim
register records
the classification
and justification
of significant
claims in
a structured
form.

A suitable
record may
contain:

\begin{enumerate}
    \item claim identifier;
    \item exact claim
    or concise
    formulation;
    \item source section
    or location;
    \item functional
    or logical
    form;
    \item epistemic
    or justificatory
    status;
    \item supporting
    premises or
    arguments;
    \item relevant
    definitions and
    formal dependencies;
    \item scope and
    conditions;
    \item unresolved
    objections or
    limitations;
    \item revision
    history, where
    relevant.
\end{enumerate}

The register
must not
replace the
arguments in
the main
text.

Its purpose
is to make
the structure
of claims
traceable and
to support
consistency checks
across the
book.

The exact
technical format
of the register
may be
determined by
the project's
registry architecture.

The present
section establishes
the conceptual
requirements
for classification
without presupposing
a particular
database schema.

\subsection{Rejected and Superseded Claims in Philosophical Revision}

The explicit
recording of
rejected and
superseded claims
has particular
importance in
a developing
philosophical system.

A revision
may correct
a conceptual
equivalence that
was previously
assumed, restrict
a claim
that was
expressed too
broadly, or
replace a
formalization
that did
not adequately
represent the
ontological
distinction
under examination.

Such revision
must distinguish
between:

\begin{itemize}
    \item the rejection
    of a claim's
    content;
    \item the rejection
    of a particular
    argument for
    that claim;
    \item the restriction
    of its scope;
    \item the replacement
    of its terminology;
    \item the reformulation
    of its formal
    representation;
    \item the supersession
    of the earlier
    construction.
\end{itemize}

These outcomes
are not equivalent.

An argument
may be rejected
while the
claim it
was intended
to support
remains defensible
through another
argument.

A definition
may be replaced
without rejecting
the phenomenon
it was intended
to describe.

A claim
may be restricted
to a narrower
domain without
being rejected
in that
domain.

The revision
record must
therefore identify
what has
changed and
why.

This is
especially important
where a
philosophical
distinction has
been developed
across several
chapters.

\subsection{Definitions and Propositions}

\paragraph{Definition 16.1: Functional or logical form.}

The functional
or logical
form of
a claim
is its
role within
the structure
of philosophical
or formal
inquiry.

\paragraph{Definition 16.2: Epistemic or justificatory status.}

The epistemic
or justificatory
status of
a claim
is its
standing according
to the
grounds on
which it
is adopted,
stipulated, argued,
derived, conditionally
supported, treated
as provisional,
retained as
a hypothesis,
left unresolved,
or rejected.

\paragraph{Definition 16.3: Two-axis claim classification.}

The two-axis
classification of
a claim
is the
ordered pair
consisting of
its functional
or logical
form and
its epistemic
or justificatory
status.

\paragraph{Definition 16.4: Claim register.}

A claim
register is
a structured
record of
significant claims,
their classifications,
their supporting
grounds, their
scope, and
their revision
history.

The following
propositions state
the governing
consequences of
the classification
system.

\paragraph{Proposition 16.1: Functional form and justificatory status are independent axes.}

The functional
or logical
form of
a claim
does not
by itself
determine its
epistemic or
justificatory
status.

\paragraph{Proposition 16.2: A formal role does not establish ontological truth.}

The classification
of a statement
as an axiom,
definition, or
formal proposition
does not
by itself
establish the
ontological truth
of its content.

\paragraph{Proposition 16.3: Derivability is relative to premises and rules.}

A formally
derived conclusion
is supported
relative to
its premises,
formal language,
and inference
rules.

\paragraph{Proposition 16.4: Justificatory status must reflect actual grounds.}

The justificatory
status assigned
to a claim
must correspond
to the
grounds actually
available for
maintaining it.

\paragraph{Proposition 16.5: Conditional support requires explicit scope.}

A claim
classified as
conditionally supported
must identify
the relevant
conditions, domain,
or interpretation
that restricts
its support.

\paragraph{Proposition 16.6: Provisional constructions must remain identifiable as provisional.}

A construction
introduced for
working purposes
must not
be silently
treated as
a foundational
commitment or
established ontological
conclusion.

\paragraph{Proposition 16.7: Open problems must remain distinguishable from proposed solutions.}

The formulation
of a possible
solution does
not, by
itself, resolve
the problem
to which
it responds.

\paragraph{Proposition 16.8: Rejection and supersession require distinct records.}

The rejection
of a claim
and the
replacement of
an earlier
formulation are
different revision
outcomes and
must be
recorded accordingly.

\paragraph{Proposition 16.9: Complex claims may require component-level classification.}

Where a
complex statement
contains logically
distinct claims
with materially
different functions
or justificatory
statuses, those
components should
be separately
classified when
necessary for
accurate assessment.

\subsection{Limits and Open Problems}

The two-axis
system provides
a structured
method for
classifying claims.
It does
not eliminate
the need
for philosophical
judgment concerning
the adequacy
of their
grounds.

Several issues
remain open
for further
development.

First, the
precise criteria
for distinguishing
a foundational
commitment from
a foundational
thesis may
require further
refinement.

Second, the
relationship between
axiomatic adoption
and philosophical
justification may
vary across
different formal
systems.

Third, the
degree of
support required
for a
claim to
move from
hypothesis to
conditionally supported
claim requires
domain-sensitive
assessment.

Fourth, the
classification of
complex arguments
may require
more detailed
records of
premise-level
statuses.

Fifth, the
project must
determine how
claim identifiers,
revision histories,
and cross-references
are represented
consistently across
its registries.

Sixth, the
relation between
philosophical
justification and
formal derivation
requires continued
attention wherever
ontological claims
are expressed
through formal
systems.

These open
problems concern
the refinement
and implementation
of the classification
method.

They do
not undermine
the central
distinction between
the two
axes.

\subsection{Concluding Principle}

A philosophical
system must
identify both
the role
of its
claims and
the grounds
on which
they are
maintained.

Functional or
logical form
describes what
a claim
does within
the system.

Epistemic or
justificatory status
describes the
standing of
the claim
according to
its available
support.

The two
dimensions must
be recorded
independently.

The governing
principle may
therefore be
stated as
follows:

\begin{quote}
No claim acquires
ontological truth
merely by being
classified as
foundational, axiomatic,
definitional, propositional,
or argumentative.
Its functional
role and
its justificatory
standing must
be identified
separately, and
its ontological
force must
be assessed
according to
the grounds
that support
its content.
\end{quote}

This distinction
provides a
methodological safeguard
for the
reconstruction of
Book I.

It allows
the system
to state
its commitments
clearly, formalize
its concepts
without overstating
their significance,
derive conclusions
without concealing
their premises,
and preserve
unresolved questions
without prematurely
converting them
into established
doctrine.

The next
section examines
the limits
of formalization
and the
relation between
formal representation
and ontological
commitment.

% =============================================================================
% END OF SECTION 16 OF CHAPTER 07
% =============================================================================